\begin{helper}[An earlier PE cycle is internal at the first later non-RE slot]
Let a certified Tashkinov series have an earlier nonterminal PE slot (r^+),
with recorded pair (\{\gamma,\delta\}), and let (C) be a
\((\gamma,\delta)\)-cycle under the coloring at that slot.  Suppose (C)
meets the ladder tree preceding the earlier PE and the terminal slot is PE.
Then there is a later slot (p^+), of kind SE or PE, such that all slots
strictly between (r^+) and (p^+) are RE, the coloring preceding (p^+)
equals the coloring at (r^+), and every exact edge identity on (C) is
internal to the ladder tree preceding (p^+).
\end{helper}
\begin{proof}
Because the earlier slot is not terminal and the terminal slot is PE, the
finite set of later slots whose kind is not RE is nonempty.  Choose its least
member (p^+).  It cannot be the initial slot, and certification by
\noderef{TashkinovSeries} and \noderef{TashkinovSeriesStep} leaves only SE,
PE, and RE at a positive slot.
Its choice excludes RE, so it is SE or PE.  Minimality says that every slot
strictly between (r^+) and (p^+) is RE.  At an RE step the
`currentColoring` field is equality with the preceding coloring.  Composing
these equalities along the finite interval identifies the coloring preceding
(p^+) with the coloring at (r^+).  Transport (C) across this equality.
Here (C) is the simple closed incidence walk supplied by
\noderef{TwoColorCycle}.

Assume, for a contradiction, that an exact cycle edge is not internal to the
tree (T_p) preceding (p^+).  The ladder trees are nested by the
`orderedExtension` fields of \noderef{TashkinovSeriesStep}; applying
\noderef{TreeSequenceInitialSegmentVertexSetSubset}, the supplied vertex of
(C) in the tree preceding (r^+) also lies in (T_p).  Traverse
the transported cycle from that vertex and choose its first crossing of
(T_p).  Its edge-node is an exact edge (f) of the cycle.  The incidence
graph definition \noderef{TwoColorIncidenceGraph} identifies the adjacent
vertex-nodes with the two endpoints of (f); hence one endpoint (x) lies in
(T_p), the other endpoint (y)
lies outside (T_p), and \noderef{BoundaryEdges} gives
(f\in\partial(T_p)).

Follow the cycle backwards from the edge-node (f) to the last visit, before
this crossing, to the earlier ladder tree.  Call that return vertex (u).
The chosen segment is a path because it is a proper segment of the simple
cycle.  Its support is contained in the cycle support.  All its vertex-nodes
lie in (T_p): otherwise the traversal would have crossed the boundary of
(T_p) before the chosen first crossing.  Moreover (u) lies in the earlier
tree by construction.

These data fill every field of
\noderef{RevisitingExtensionCertificate} for the stage preceding (p^+):
take earlier slot (r^+), its positive index and PE kind, the recorded colors
(\gamma,\delta), the transported cycle, the given earlier-tree meeting,
the outside vertex (y), boundary edge (f), endpoints (x,y), return
vertex (u), and the internal return path just constructed.  The interval of
RE slots verifies `interveningRevisiting`.  But an SE or PE certificate at
(p^+) contains `noRevisiting`, which denies the existence of precisely this
certificate.  The contradiction shows that every edge of (C) is internal
to (T_p), in the precise sense of \noderef{InternalEdges}.  Together with
the composed RE coloring equality, this proves all asserted conclusions.
\end{proof}
